\paragrafos{Τυποποιημένη μορφή μεγάλων και μικρών αριθμών}{Α.7.10, «Τυποποιημένη μορφή μεγάλων και μικρών αριθμών»}

\begin{anagke}
Η Α3 επισκέπτεται το εργαστήριο βιολογίας του σχολείου και η κυρία Αλεξάνδρα δείχνει ένα μικροσκόπιο. Κάτω από τον φακό υπάρχουν μια τρίχα, ένα ερυθρό αιμοσφαίριο και ένα βακτήριο. Η καθηγήτρια γράφει στον πίνακα τα μεγέθη: η τρίχα έχει πάχος $0{,}000\,07$ μ., το αιμοσφαίριο $0{,}000\,007$ μ., το βακτήριο $0{,}000\,002$ μ. και ένας ιός $0{,}000\,000\,1$ μ. Ο Στέλιος προσπαθεί να τα συγκρίνει και χάνεται στα μηδενικά: «Ποιο είναι μικρότερο, το $0{,}000\,007$ ή το $0{,}000\,002$;».

Η Πολυξένη θυμάται ότι για τους πολύ μεγάλους αριθμούς η κυρία Αλεξάνδρα τους είχε δείξει να τους γράφουν με δυνάμεις του $10$: $1{,}496\cdot10^{8}$ χιλιόμετρα η απόσταση Γης -- Ήλιου. «Μπορούμε να κάνουμε το ίδιο και για τους πολύ μικρούς;» ρωτάει. «Τώρα που ξέρουμε τους αρνητικούς εκθέτες, ναι», απαντάει η καθηγήτρια. Το $0{,}000\,07$ γράφεται $7\cdot10^{-5}$, και ο ιός $1\cdot10^{-7}$.

Η Μυρτώ προσθέτει μια ιστορική σημείωση: «Αυτά τα μεγέθη μπορούν να τα μετρήσουν χάρη στα επιστημονικά όργανα και στη γλώσσα των αριθμών. Η Υπατία, η μεγάλη μαθηματικός και αστρονόμος της Αλεξάνδρειας (περίπου 370--415 μ.Χ.), είχε δουλέψει σε όργανα μέτρησης του ουρανού, τον αστρολάβο». Η τυποποιημένη μορφή φέρνει τους πολύ μεγάλους και τους πολύ μικρούς αριθμούς στο ίδιο τετράδιο, με την ίδια γλώσσα.
\end{anagke}

\begin{lexi}{τυποποιημένη μορφή, τάξη μεγέθους}{τυποποιημένη μορφή, τάξη μεγέθους}
Όπως είδαμε, \lexiword{τυποποιημένη} είναι η μορφή που έχει το ίδιο «τύπο» για όλους. Η \lexiword{τάξη μεγέθους} ενός αριθμού είναι ο εκθέτης του $10$ στην τυποποιημένη του μορφή: μας λέει «πόσο μεγάλος ή πόσο μικρός» είναι, με μια ματιά.
\end{lexi}

\begin{metafrasi}
Το πάχος της τρίχας. & $0{,}000\,07=7\cdot10^{-5}$ μ.\\
Το αιμοσφαίριο. & $0{,}000\,007=7\cdot10^{-6}$ μ.\\
Ο ιός. & $0{,}000\,000\,1=1\cdot10^{-7}$ μ.\\
Η ταχύτητα του φωτός. & $300\,000\,000=3\cdot10^{8}$ μ./s\\
Η οθόνη του κομπιουτεράκι. & $4.5\mathrm{E}{-}4=4{,}5\cdot10^{-4}$\\
\end{metafrasi}

\tmima{Η θεωρία}

\begin{orismos}{τυποποιημένη μορφή}
Κάθε θετικός ρητός μπορεί να γραφεί στη μορφή $\alpha\cdot10^{\nu}$ με $1\le\alpha<10$ και $\nu$ \textbf{ακέραιο}. Αυτή λέγεται \textbf{τυποποιημένη μορφή}. Για τους μεγάλους αριθμούς ο $\nu$ είναι θετικός, για τους μικρούς ($<1$) αρνητικός.
\end{orismos}

\begin{kanonas}{τυποποίηση}
\begin{itemize}[leftmargin=1.2em,itemsep=0.3ex,topsep=2pt]
\item Μεταφέρω την υποδιαστολή ώστε να μείνει ένα μη μηδενικό ψηφίο πριν από αυτήν και βρίσκω τον $\alpha$.
\item Ο εκθέτης $\nu$ είναι το πλήθος των θέσεων που μετακίνησα, με πρόσημο $+$ αν μετακίνησα την υποδιαστολή αριστερά (αριθμός μεγάλος) και $-$ αν τη μετακίνησα δεξιά (αριθμός μικρός).
\end{itemize}
\end{kanonas}

Έτσι $0{,}00045$: μετακινώ $4$ θέσεις δεξιά και παίρνω $4{,}5$: $0{,}00045=4{,}5\cdot10^{-4}$. Και $3\,420\,000=3{,}42\cdot10^{6}$. Για να γυρίσω πίσω, πολλαπλασιάζω τον $\alpha$ με το $10^\nu$: $3{,}2\cdot10^{-5}=0{,}000032$.

\tmimasmall{Σύγκριση και πράξεις}

Συγκρίνω πρώτα τους εκθέτες: μεγαλύτερος εκθέτης σημαίνει μεγαλύτερο αριθμό ($4\cdot10^{-5}>7\cdot10^{-6}$ αφού $-5>-6$). Αν οι εκθέτες είναι ίσοι συγκρίνω τους $\alpha$. Για πράξεις: στον πολλαπλασιασμό πολλαπλασιάζω τους $\alpha$ και προσθέτω τους εκθέτες, στη διαίρεση διαιρώ τους $\alpha$ και αφαιρώ τους εκθέτες, και μετά ξαναγράφω σε τυποποιημένη μορφή αν χρειάζεται:
$$(3\cdot10^{-3})(2\cdot10^{5})=6\cdot10^{2},\qquad(6\cdot10^{8}):(2\cdot10^{-2})=3\cdot10^{10},\qquad(5\cdot10^{4})(4\cdot10^{-2})=20\cdot10^{2}=2\cdot10^{3}.$$

\begin{prosoxi}
\begin{itemize}[leftmargin=1.2em,itemsep=0.3ex,topsep=2pt]
\item Ο $\alpha$ είναι πάντα από $1$ έως λιγότερο από $10$. Το $12\cdot10^{-4}$ δεν είναι τυποποιημένο: είναι $1{,}2\cdot10^{-3}$.
\item Στους μικρούς αριθμούς ο εκθέτης είναι \emph{αρνητικός}: το $0{,}00045$ δεν είναι $4{,}5\cdot10^{4}$ αλλά $4{,}5\cdot10^{-4}$.
\item Ο εκθέτης $-4$ δεν σημαίνει «τέσσερα μηδενικά» αλλά «η υποδιαστολή πήγε τέσσερις θέσεις δεξιά»: μετά την υποδιαστολή υπάρχουν $3$ μηδενικά και μετά το $4$.
\item Στη σύγκριση: $10^{-5}$ είναι μεγαλύτερο από $10^{-6}$. Όσο πιο αρνητικός ο εκθέτης, τόσο μικρότερος ο αριθμός.
\end{itemize}
\end{prosoxi}

\tmima{Λυμένα παραδείγματα}

\paradeigma{Μικροί αριθμοί}
\tixero Οι $0{,}00045$, $0{,}000\,000\,7$, $0{,}0302$.\par
\tithelo Να τους γράψω σε τυποποιημένη μορφή.\par
\possindeo Μετακινώ την υποδιαστολή δεξιά μέχρι ένα μη μηδενικό ψηφίο πριν από αυτήν. $0{,}00045$: $4$ θέσεις, $4{,}5\cdot10^{-4}$. $0{,}000\,000\,7$: $7$ θέσεις, $7\cdot10^{-7}$. $0{,}0302$: $2$ θέσεις, $3{,}02\cdot10^{-2}$.\par
\apantisi $4{,}5\cdot10^{-4}$, $7\cdot10^{-7}$, $3{,}02\cdot10^{-2}$.\par
\elegxos $4{,}5\cdot10^{-4}=\frac{4{,}5}{10\,000}=0{,}00045$.

\paradeigma{Από τυποποιημένη σε δεκαδική}
\tixero Οι $3{,}2\cdot10^{-5}$, $9{,}1\cdot10^{-3}$, $2{,}5\cdot10^{6}$.\par
\tithelo Να τους γράψω κανονικά.\par
\possindeo Αρνητικός εκθέτης: υποδιαστολή αριστερά. $3{,}2\cdot10^{-5}=0{,}000032$. $9{,}1\cdot10^{-3}=0{,}0091$. Θετικός εκθέτης: υποδιαστολή δεξιά. $2{,}5\cdot10^6=2\,500\,000$.\par
\apantisi $0{,}000032$, $0{,}0091$, $2\,500\,000$.\par
\elegxos $3{,}2\cdot10^{-5}$: $5$ θέσεις αριστερά από το $3$: $0{,}000032$.

\paradeigma{Σύγκριση μεγεθών}
\tixero Η τρίχα $7\cdot10^{-5}$ μ., το αιμοσφαίριο $7\cdot10^{-6}$ μ., το βακτήριο $2\cdot10^{-6}$ μ., ο ιός $1\cdot10^{-7}$ μ.\par
\tithelo Να τα βάλω σε σειρά από το μικρότερο και να βρω πόσες φορές η τρίχα είναι χοντρότερη από το αιμοσφαίριο.\par
\possindeo Συγκρίνω εκθέτες: $-7<-6<-5$. Ο ιός ο μικρότερος. Μετά τα δύο με εκθέτη $-6$: $2<7$, δηλαδή το βακτήριο και μετά το αιμοσφαίριο. Τελευταία η τρίχα. Ο λόγος: $\frac{7\cdot10^{-5}}{7\cdot10^{-6}}=10^{-5-(-6)}=10$.\par
\apantisi $1\cdot10^{-7}<2\cdot10^{-6}<7\cdot10^{-6}<7\cdot10^{-5}$. Η τρίχα είναι $10$ φορές χοντρότερη.\par
\elegxos $0{,}00007=10\cdot0{,}000007$.

\paradeigma{Πράξεις}
\tixero Οι $(3\cdot10^{-3})\cdot(2\cdot10^{5})$, $(6\cdot10^{8}):(2\cdot10^{-2})$, $(5\cdot10^{4})\cdot(4\cdot10^{-2})$.\par
\tithelo Να τις υπολογίσω σε τυποποιημένη μορφή.\par
\possindeo $3\cdot2=6$ και $-3+5=2$: $6\cdot10^2$. $6:2=3$ και $8-(-2)=10$: $3\cdot10^{10}$. $5\cdot4=20$ και $4-2=2$: $20\cdot10^2=2\cdot10^3$.\par
\apantisi $6\cdot10^{2}$, $3\cdot10^{10}$, $2\cdot10^{3}$.\par
\elegxos Το τρίτο: $50\,000\cdot0{,}04=2\,000$.

\paradeigma{Το φως του Ήλιου}
\tixero Το φως ταξιδεύει με $3\cdot10^{8}$ μ./s και ο Ήλιος απέχει $1{,}5\cdot10^{11}$ μ.\par
\tithelo Σε πόσο χρόνο φτάνει το φως στη Γη.\par
\possindeo $\frac{1{,}5\cdot10^{11}}{3\cdot10^{8}}=0{,}5\cdot10^{3}=5\cdot10^{2}$ s.\par
\apantisi $500$ s, περίπου $8{,}3$ λεπτά.\par
\elegxos $500\cdot3\cdot10^8=1{,}5\cdot10^{11}$.

\tmima{Ασκήσεις}

\zoni{Το κατάλαβα}{μικροί και μεγάλοι αριθμοί}

\ex Γράψε σε τυποποιημένη μορφή: α) $0{,}0005$ \quad β) $0{,}000\,072$ \quad γ) $0{,}3$ \quad δ) $45\,000$ \quad ε) $0{,}000\,000\,8$ \quad στ) $6\,300\,000$.

\ex Γράψε σε δεκαδική μορφή: α) $3{,}2\cdot10^{-5}$ \quad β) $9{,}1\cdot10^{-3}$ \quad γ) $4\cdot10^{-1}$ \quad δ) $2{,}5\cdot10^{6}$ \quad ε) $1{,}01\cdot10^{-4}$ \quad στ) $7\cdot10^{0}$.

\ex Γράψε σωστά σε τυποποιημένη μορφή: α) $0{,}3\cdot10^{-2}$ \quad β) $12\cdot10^{-4}$ \quad γ) $5{,}5\cdot10^{-3}$ \quad δ) $10\cdot10^{2}$.

\ex Σύγκρινε: α) $4\cdot10^{-5}$ και $7\cdot10^{-6}$ \quad β) $2{,}5\cdot10^{-3}$ και $2{,}5\cdot10^{-2}$ \quad γ) $6\cdot10^{4}$ και $9\cdot10^{-4}$ \quad δ) $3{,}1\cdot10^{-8}$ και $3{,}09\cdot10^{-8}$.

\ex Γράψε σε αύξουσα σειρά: $3\cdot10^{-3}$, $5\cdot10^{-5}$, $2\cdot10^{2}$, $9\cdot10^{-1}$, $1{,}2\cdot10^{-3}$.

\ex Υπολόγισε σε τυποποιημένη μορφή: α) $(3\cdot10^{-3})(2\cdot10^{5})$ \quad β) $(6\cdot10^{8}):(2\cdot10^{-2})$ \quad γ) $(5\cdot10^{4})(4\cdot10^{-2})$ \quad δ) $(8\cdot10^{-6}):(4\cdot10^{-3})$.

\ex Ο υπολογιστής δείχνει $4.5\mathrm{E}{-}4$, $2.5\mathrm{E}7$ και $7\mathrm{E}{-}9$. Ποιοι αριθμοί είναι;

\ex Σωστό ή λάθος; α) $0{,}004=4\cdot10^{-3}$. \quad β) Το $12\cdot10^{-4}$ είναι τυποποιημένο. \quad γ) Το $10^{-5}$ είναι μικρότερο από το $10^{-6}$. \quad δ) $2\cdot10^{-3}=0{,}002$. \quad ε) Όσο μεγαλύτερος ο εκθέτης, τόσο μεγαλύτερος ο αριθμός. \quad στ) $5{,}5\cdot10^{-3}>5{,}5\cdot10^{-2}$.

\zoni{Το χρησιμοποιώ}{ο μικρόκοσμος και ο μακρόκοσμος}

\ex \textbf{Ο μικρόκοσμος.} Τρίχα $7\cdot10^{-5}$ μ., αιμοσφαίριο $7\cdot10^{-6}$ μ., βακτήριο $2\cdot10^{-6}$ μ., ιός $10^{-7}$ μ. Βάλε τα σε αύξουσα σειρά και βρες πόσες φορές η τρίχα είναι χοντρότερη από το αιμοσφαίριο.

\ex \textbf{Το ηλεκτρόνιο.} Η μάζα του ηλεκτρονίου είναι $9{,}1\cdot10^{-31}$ kg. Πόσα μηδενικά υπάρχουν μετά την υποδιαστολή πριν από το $9$ αν τη γράψεις δεκαδικά;

\ex \textbf{Το φως.} Το φως ταξιδεύει με $3\cdot10^{8}$ μ./s και ο Ήλιος απέχει από τη Γη $1{,}5\cdot10^{11}$ μ. Σε πόσο χρόνο φτάνει το φως στη Γη;

\ex \textbf{Το DNA.} Το πλάτος του DNA είναι $2$ nm και $1\ \mathrm{nm}=10^{-9}$ μ. Γράψε το σε μέτρα σε τυποποιημένη μορφή.

\ex \textbf{Το χαρτί.} Ένα φύλλο χαρτιού έχει πάχος $10^{-4}$ μ. Πόσο είναι το πάχος $500$ φύλλων σε μέτρα και σε εκατοστά;

\ex \textbf{Η άμμος.} Ένας κόκκος άμμου έχει μάζα $10^{-6}$ kg. Πόσοι κόκκοι ζυγίζουν $1$ kg;

\zoni{Το σκέφτομαι}{ο εκθέτης που μικραίνει}

\ex Γιατί ο εκθέτης είναι αρνητικός για αριθμούς μικρότερους από το $1$;

\ex Ποια είναι η τυποποιημένη μορφή του $1$; Και του $10$;

\ex Σύγκρινε τα $2\cdot10^{-3}$ και $5\cdot10^{-4}$ με δύο τρόπους (δεκαδικά και εκθέτες).

\ex Αν πολλαπλασιάσεις το $3\cdot10^{-4}$ με το $10^{6}$ τι παίρνεις;

\ex Γιατί το $12\cdot10^{-4}$ δεν είναι τυποποιημένο; Γράψε το σωστά.

\zoni{Για το σπίτι}{ψάξε τα μεγέθη}

\ex Βρες το μέγεθος ενός ατόμου, ενός ιού και ενός βακτηρίου από αξιόπιστη πηγή και γράψε τα σε τυποποιημένη μορφή και σε αύξουσα σειρά.
